% This contents of this file will be inserted into the _Solutions version of the
% output tex document.  Here's an example:

% If assignment with subquestion (1.a) requires a written response, you will
% find the following flag within this document: <SCPD_SUBMISSION_TAG>_1a
% In this example, you would insert the LaTeX for your solution to (1.a) between
% the <SCPD_SUBMISSION_TAG>_1a flags.  If you also constrain your answer between the
% START_CODE_HERE and END_CODE_HERE flags, your LaTeX will be styled as a
% solution within the final document.

% Please do not use the '<SCPD_SUBMISSION_TAG>' character anywhere within your code.  As expected,
% that will confuse the regular expressions we use to identify your solution.

\def\assignmentnum{3 }
\def\assignmentname{Route Planning}
\def\assignmenttitle{XCS221 Assignment \assignmentnum --- \assignmentname}
\input{macros}
\begin{document}
\pagestyle{myheadings} \markboth{}{\assignmenttitle}

% <SCPD_SUBMISSION_TAG>_entire_submission

This handout includes space for every question that requires a written response.
Please feel free to use it to handwrite your solutions (legibly, please).  If
you choose to typeset your solutions, the |README.md| for this assignment includes
instructions to regenerate this handout with your typeset \LaTeX{} solutions.
\ruleskip

\LARGE
1.a
\normalsize

% <SCPD_SUBMISSION_TAG>_1a
\begin{answer}
  % ### START CODE HERE ###
  % Kept for reference, not rendered. How I got here:
  %  - Cost((x, y), a) is priced on the state being LEFT, and depends only on
  %    x: 1 + x for x >= 0, and 1 for x < 0. Direction does not matter, so a
  %    north step from x = 2 costs 3, not 1.
  %  - each action moves one coordinate by one, so every step is priced
  %    separately; the formula is applied per step, never once per leg
  %  - checked by hand: (2, 1) north-first costs 1 + 1 + 2 = 4, while
  %    east-first costs 1 + 2 + 3 = 6; (3, 0) costs 1 + 2 + 3 = 6
  %  - guesses that failed a hand check: 1 + m, 1 + m + n, 2n + m, 2m + n
  %  - closed form of 1 + 2 + ... + m by writing the series forwards and
  %    backwards: m columns each summing to m + 1, two copies of the series,
  %    so the sum is m(m + 1)/2. Works for odd m, where pairing leaves a middle
  %    term.
  \textbf{Minimum cost.}
  \[
    \frac{m(m+1)}{2} + n .
  \]

  \medskip
  Every step is priced by the $x$ coordinate of the location being left, whatever the direction.
  The $n$ north steps each cost exactly 1 if they are taken while $x = 0$, contributing $n$. The $m$
  east steps have to leave from $x = 0, 1, \dots, m-1$ in turn, costing $1, 2, \dots, m$, which sums
  to $m(m+1)/2$.

  \medskip
  \textbf{Path.} Take all $n$ steps north from the origin first, then all $m$ steps east.

  \medskip
  \textbf{Uniqueness.} The path is unique. Any north step taken after moving east costs more than 1,
  since it leaves from $x > 0$, and a detour into negative $x$ adds at least one step west and one
  step back east, each costing at least 1, without making any other step cheaper.
  % ### END CODE HERE ###
\end{answer}
% <SCPD_SUBMISSION_TAG>_1a
\clearpage

\LARGE
1.b
\normalsize

% <SCPD_SUBMISSION_TAG>_1b
\begin{answer}
  % ### START CODE HERE ###
  % ### END CODE HERE ###
\end{answer}
% <SCPD_SUBMISSION_TAG>_1b
\clearpage

\LARGE
1.c
\normalsize

% <SCPD_SUBMISSION_TAG>_1c
\begin{answer}
  % ### START CODE HERE ###
  % ### END CODE HERE ###
\end{answer}
% <SCPD_SUBMISSION_TAG>_1c
\clearpage

\LARGE
2.b
\normalsize

% <SCPD_SUBMISSION_TAG>_2b
\begin{answer}
  % ### START CODE HERE ###
  % ### END CODE HERE ###
\end{answer}
% <SCPD_SUBMISSION_TAG>_2b


\LARGE
2.c
\normalsize

% <SCPD_SUBMISSION_TAG>_2c
\begin{answer}
  % ### START CODE HERE ###
  % ### END CODE HERE ###
\end{answer}
% <SCPD_SUBMISSION_TAG>_2c

\clearpage

\LARGE
3.b
\normalsize

% <SCPD_SUBMISSION_TAG>_3b
\begin{answer}
  % ### START CODE HERE ###
  % ### END CODE HERE ###
\end{answer}
% <SCPD_SUBMISSION_TAG>_3b
\clearpage

\LARGE
3.c
\normalsize

% <SCPD_SUBMISSION_TAG>_3c
\begin{answer}
  % ### START CODE HERE ###
  % ### END CODE HERE ###
\end{answer}
% <SCPD_SUBMISSION_TAG>_3c

\LARGE
3.d
\normalsize

% <SCPD_SUBMISSION_TAG>_3d
\begin{answer}
  % ### START CODE HERE ###
  % ### END CODE HERE ###
\end{answer}
% <SCPD_SUBMISSION_TAG>_3d
\clearpage

\LARGE
5.a
\normalsize

% <SCPD_SUBMISSION_TAG>_5a
\begin{answer}
  % ### START CODE HERE ###
  % ### END CODE HERE ###
\end{answer}
% <SCPD_SUBMISSION_TAG>_5a

\LARGE
5.b
\normalsize

% <SCPD_SUBMISSION_TAG>_5b
\begin{answer}
  % ### START CODE HERE ###
  % ### END CODE HERE ###
\end{answer}
% <SCPD_SUBMISSION_TAG>_5b

\clearpage

\LARGE
6.a
\normalsize

% <SCPD_SUBMISSION_TAG>_6a
\begin{answer}
  % ### START CODE HERE ###
  % ### END CODE HERE ###
\end{answer}
% <SCPD_SUBMISSION_TAG>_6a

\LARGE
6.b
\normalsize

% <SCPD_SUBMISSION_TAG>_6b
\begin{answer}
  % ### START CODE HERE ###
  % ### END CODE HERE ###
\end{answer}
% <SCPD_SUBMISSION_TAG>_6b

\LARGE
6.c
\normalsize

% <SCPD_SUBMISSION_TAG>_6c
\begin{answer}
  % ### START CODE HERE ###
  % ### END CODE HERE ###
\end{answer}
% <SCPD_SUBMISSION_TAG>_6c

\LARGE
6.d
\normalsize

% <SCPD_SUBMISSION_TAG>_6d
\begin{answer}
  % ### START CODE HERE ###
  % ### END CODE HERE ###
\end{answer}
% <SCPD_SUBMISSION_TAG>_6d
\clearpage

% <SCPD_SUBMISSION_TAG>_entire_submission

\end{document}